\chapter{Normalisierte ECS-Datenmodelle und domänenspezifische Handles zur statischen Konfliktanalyse von Systemen in einer C++-Game-Engine}\label{ch:vorschlag1}

\section{Forschungsfrage}

\begin{quote}
    Wie kann eine normalisierte relationale Modellierung von ECS-Datenstrukturen genutzt werden, um Lese- und Schreibzugriffe von ECS-Systemen statisch zu beschreiben, und inwieweit tragen domänenspezifische Handles dazu bei, Konflikte zwischen Systemen zu erkennen oder auszuschließen?
\end{quote}


\section{Methodik}

Die Arbeit untersucht, ob domänenspezifische Handles im Zusammenhang mit einem relational modellierten ECS-Datenmodell dabei unterstützen, Scheduling-Konflikte auf Komponentenebene statisch zu analysieren.\\

Hierzu wird zunächst \textit{ECS Core} nach Redmond et {al.} als theoretischer Referenzrahmen herangezogen.
Die Arbeit stellt Gemeinsamkeiten und Abgrenzungen zwischen dem dort beschriebenen Modell und der Implementierung in \texttt{helios::ecs} heraus.
Dies betrifft insbesondere das Query-Modell, die Beschreibung von Komponentenzugriffen sowie die Einordnung serieller und nebenläufiger Systemausführung.\\

Der theoretische Teil motiviert anhand der Normalformtheorie im relationalen Modell eine entitätszentrische Relation

\[
    K(e, k_1, k_2, \ldots, k_n)
\]

\noindent
sowie eine komponentenzentrische Zerlegung

\[
    K_1(e, k_1),\quad K_2(e, k_2),\quad \ldots,\quad K_n(e, k_n)
\]

\noindent
von Entitäten und Komponenten.
Die Verwendung von \textit{Sparse Sets} in \texttt{helios::ecs} wird dabei als konkrete Speicherstruktur für die komponentenzentrische Datenhaltung betrachtet.\\

Im Anschluss werden ECS-Systeme über ihre Lese- und Schreibzugriffe auf Komponentenrelationen beschrieben.\\
Dabei sollen statische Informationen zur Ableitung möglicher Konflikte beim Scheduling genutzt werden, wobei die Tragfähigkeit dieses Ansatzes im Verlauf der Arbeit geprüft wird: In welchem Umfang erlauben domänenspezifische Handles eine Einordnung möglicher Konflikte auf Komponentenebene?\\

Die Untersuchung erfolgt in vier Schritten:

\begin{enumerate}
    \item \textbf{Modellbildung:} Entitäten, Komponenten und Komponentenrelationen werden relational beschrieben und hinsichtlich ihrer Normalformeigenschaften eingeordnet.

    \item \textbf{Abbildung auf \texttt{helios::ecs}:} Die theoretische Zerlegung wird auf Komponenten, Queries und Sparse Sets in \texttt{helios::ecs} bezogen.

    \item \textbf{Zugriffsmodell:} ECS-Systeme werden über Lese- und Schreibmengen auf Komponentenrelationen beschrieben.

    \item \textbf{Konfliktanalyse:} Für ein Referenzszenario wird geprüft, ob sich aus den Zugriffsmengen und den domänenspezifischen Handles Aussagen über konfliktarme Systemausführung ableiten lassen.
    Ergänzend wird untersucht, ob sich für das betrachtete System eine Nähe zur Konfliktserialisierbarkeit formal beschreiben lässt.
\end{enumerate}


\section{Vorläufige Gliederung}

\begin{enumerate}
    \item Einleitung und Zielsetzung
    \item Grundlagen: ECS, relationales Modell, komponentenzentrische Speicherung und Systeme
    \item \textit{ECS Core}-Modell und Abgrenzung zu \texttt{helios::ecs}
    \item Relationales ECS-Modell und Normalformanalyse
    \item Abbildung auf \texttt{helios::ecs}: Komponenten und Queries
    \item Systemzugriffe als Lese-/Schreibmengen
    \item Domänenspezifische Handles als statische Systemgrenzen
    \item Typisierte Systemsignaturen: Lese-/Schreibmengen und Domänen
    \item Konfliktanalyse von ECS-Systemen
    \item Referenzszenario: Kollisions- und Partikelsysteme
    \item Evaluation: Ableitung von Konflikten im Referenzszenario
    \item Diskussion
    \item Fazit
\end{enumerate}